% math0 manuscript. Build with: latexmk -pdf paper.tex
\documentclass[11pt]{amsart}
\usepackage[margin=1.1in]{geometry}
\usepackage{amsmath,amssymb,amsthm,mathtools}
\usepackage{url}
\usepackage{hyperref}
\usepackage{xcolor}
\usepackage{listings}
\usepackage{newunicodechar}
\usepackage{booktabs}

\lstdefinelanguage{lean}{
  keywords={theorem,lemma,def,abbrev,example,fun,by,have,show,from,let,in,if,then,else,match,with,
    namespace,end,open,import,section,variable,instance,structure,where,calc,obtain,intro,exact,
    refine,rcases,rintro,simp,omega,decide,norm_num,linarith,nlinarith,positivity,induction,cases,
    noncomputable,private},
  sensitive=true, comment=[l]{--}, morecomment=[s]{/-}{-/}, string=[b]"}
\lstset{language=lean, basicstyle=\ttfamily\footnotesize, columns=fullflexible, keepspaces=true,
  breaklines=true, keywordstyle=\color{blue!60!black}\bfseries, commentstyle=\color{gray},
  frame=single, extendedchars=true, inputencoding=utf8,
  literate={ℕ}{{$\mathbb{N}$}}1 {ℝ}{{$\mathbb{R}$}}1 {ℤ}{{$\mathbb{Z}$}}1 {∑}{{$\sum$}}1
    {∈}{{$\in$}}1 {√}{{$\surd$}}1 {π}{{$\pi$}}1 {≤}{{$\le$}}1 {∧}{{$\wedge$}}1
    {∨}{{$\vee$}}1 {→}{{$\to$}}1 {∀}{{$\forall$}}1}
% Unicode commonly found in Lean sources.
\newunicodechar{ℕ}{\ensuremath{\mathbb{N}}}
\newunicodechar{ℤ}{\ensuremath{\mathbb{Z}}}
\newunicodechar{ℚ}{\ensuremath{\mathbb{Q}}}
\newunicodechar{ℝ}{\ensuremath{\mathbb{R}}}
\newunicodechar{ℂ}{\ensuremath{\mathbb{C}}}
\newunicodechar{∀}{\ensuremath{\forall}}
\newunicodechar{∃}{\ensuremath{\exists}}
\newunicodechar{→}{\ensuremath{\to}}
\newunicodechar{←}{\ensuremath{\leftarrow}}
\newunicodechar{↔}{\ensuremath{\leftrightarrow}}
\newunicodechar{↦}{\ensuremath{\mapsto}}
\newunicodechar{¬}{\ensuremath{\neg}}
\newunicodechar{∧}{\ensuremath{\wedge}}
\newunicodechar{∨}{\ensuremath{\vee}}
\newunicodechar{≤}{\ensuremath{\le}}
\newunicodechar{≥}{\ensuremath{\ge}}
\newunicodechar{≠}{\ensuremath{\ne}}
\newunicodechar{∈}{\ensuremath{\in}}
\newunicodechar{∉}{\ensuremath{\notin}}
\newunicodechar{⊆}{\ensuremath{\subseteq}}
\newunicodechar{∩}{\ensuremath{\cap}}
\newunicodechar{∪}{\ensuremath{\cup}}
\newunicodechar{∑}{\ensuremath{\sum}}
\newunicodechar{∏}{\ensuremath{\prod}}
\newunicodechar{⟨}{\ensuremath{\langle}}
\newunicodechar{⟩}{\ensuremath{\rangle}}
\newunicodechar{·}{\ensuremath{\cdot}}
\newunicodechar{∘}{\ensuremath{\circ}}
\newunicodechar{⁻}{\ensuremath{^{-}}}
\newunicodechar{¹}{\ensuremath{^{1}}}
\newunicodechar{²}{\ensuremath{^{2}}}
\newunicodechar{₀}{\ensuremath{_{0}}}
\newunicodechar{₁}{\ensuremath{_{1}}}
\newunicodechar{₂}{\ensuremath{_{2}}}
\newunicodechar{α}{\ensuremath{\alpha}}
\newunicodechar{β}{\ensuremath{\beta}}
\newunicodechar{ε}{\ensuremath{\varepsilon}}
\newunicodechar{δ}{\ensuremath{\delta}}
\newunicodechar{π}{\ensuremath{\pi}}
\newunicodechar{σ}{\ensuremath{\sigma}}
\newunicodechar{φ}{\ensuremath{\varphi}}
\newunicodechar{∞}{\ensuremath{\infty}}
\newunicodechar{⌊}{\ensuremath{\lfloor}}
\newunicodechar{⌋}{\ensuremath{\rfloor}}
\newunicodechar{⌈}{\ensuremath{\lceil}}
\newunicodechar{⌉}{\ensuremath{\rceil}}
\newunicodechar{∣}{\ensuremath{\mid}}
\newunicodechar{≡}{\ensuremath{\equiv}}
\newunicodechar{×}{\ensuremath{\times}}
\newunicodechar{⧸}{\ensuremath{/}}
\newunicodechar{‖}{\ensuremath{\|}}
\newunicodechar{√}{\ensuremath{\surd}}
\newunicodechar{∅}{\ensuremath{\emptyset}}
\newunicodechar{⦃}{\ensuremath{\{\!|}}
\newunicodechar{⦄}{\ensuremath{|\!\}}}
\newunicodechar{𝓝}{\ensuremath{\mathcal{N}}}
\newunicodechar{≫}{\ensuremath{\gg}}
\newunicodechar{≪}{\ensuremath{\ll}}

\theoremstyle{plain}
\newtheorem{theorem}{Theorem}
\newtheorem{lemma}[theorem]{Lemma}
\newtheorem{proposition}[theorem]{Proposition}
\newtheorem{corollary}[theorem]{Corollary}
\newtheorem{conjecture}[theorem]{Conjecture}
\theoremstyle{definition}
\newtheorem{definition}[theorem]{Definition}
\newtheorem{remark}[theorem]{Remark}

\emergencystretch=3em
\raggedbottom
% Lean identifiers in running text (breakable at dots and underscores).
\expandafter\def\expandafter\UrlBreaks\expandafter{\UrlBreaks\do\_\do\.}
\DeclareUrlCommand\decl{\urlstyle{tt}}
\newcommand{\N}{\mathbb{N}}

\title[Turns of the spiral of Theodorus]{The numbers of complete triangles in consecutive turns of the spiral of Theodorus differ by 18, 19, 20 or 21}
\author{math0}
\address{Produced by Claude (an AI model developed by Anthropic) in the math0 project}
\date{October 24, 2026}

\begin{document}
\begin{abstract}
In the spiral of Theodorus the $i$-th right triangle has legs $\sqrt i$ and $1$ and angle $\arctan(1/\sqrt i)$ at the centre; let $c(n)$ be the number of triangles that lie entirely in the $n$-th turn.
We prove that $c(n)-c(n-1)\in\{18,19,20,21\}$ for every $n\ge2$, which is the conjecture stated by S\'ebastien Dumortier in 2010 in the OEIS entry A172164.
Two elementary one-step inequalities give $|c(n)-c(n-1)-2\pi^2|<2.13$ for $n\ge71$, and a rigorous interval computation covers $n\le70$.
We found no earlier proof and classify the result as new, with the caveat that the argument is an effective form of Hlawka's classical asymptotics.
The proof is verified in Lean~4 with Mathlib against the frozen, fidelity-reviewed statement that math0 formalized from the OEIS entry, whose definitions are checked by the kernel against all $69$ listed terms.
\end{abstract}
\maketitle

\section{Introduction}

The spiral of Theodorus (\emph{Pythagoras's snail} in the OEIS) is made of right triangles around a common vertex: the first has legs $1,1$, and each next one has as legs the previous hypotenuse and a segment of length $1$.
So triangle $i\ge1$ has legs $\sqrt i$ and $1$ and angle $\arctan(1/\sqrt i)$ at the centre, and the first $k$ triangles fill the angle
\[
  S(k)=\sum_{i=1}^{k}\arctan\frac1{\sqrt i}\qquad(S(0)=0).
\]
Triangle $i$ occupies the angular sector $[S(i-1),S(i)]$, and the $n$-th turn is the range of angles $[2\pi(n-1),2\pi n]$.

A072895 \cite{A072895} (R.~G.~Wilson~v, 2002) lists the least $k$ with $S(k)>2\pi n$, and A137515 \cite{A137515} (S.~Dumortier, 2008) lists the number $K_n=\text{A072895}(n)-1$ of triangles in the first $n$ turns: $16,53,109,185,\dots$
A comment in A137515 says that ``from one term to the next one, the number added grows by 18, 19, 20 or 21 (tested up to 5000 terms)'' (restated by H.~P.~Dale, 2019).
Finally A172164 \cite{A172164} (S.~Dumortier, 2010) lists $a(n)=c(n)-c(n-1)$, where $c(n)$ is the number of triangles entirely contained in turn $n$, and conjectures that $a(n)\in\{18,19,20,21\}$ for all $n\ge2$ (Conjecture~\ref{conj}).
R.~G.~Wilson~v and H.~Kociemba added frequency counts of the four values over the first $10^5$ and $10^7$ turns (2013).

The conjecture is explained by the asymptotics of $S$.
Hlawka \cite{Hl80} showed that $S(k)=2\sqrt k+\kappa+o(1)$, where $\kappa=-2.1577829966\ldots$ is Hlawka's \emph{Schneckenkonstante} (denoted $K$ in A105459 \cite{A105459}); see Davis \cite[pp.~157--167]{Da93} for an English account and Brink \cite{Br12} for the relation of the spiral to zeta values at the half-integers.
Hence $K_n\approx(\pi n-\kappa/2)^2$; Kociemba observed numerically (A072895) that $\lfloor(\pi n-\kappa/2)^2-1/6\rfloor=\text{A072895}(n)$ for $n<10^9$.
Since the second difference $\Delta^2$ of $(\pi n+\text{const})^2$ is $2\pi^2=19.739\ldots$ and $|\Delta^2\lfloor y_n\rfloor-\Delta^2y_n|<2$ for real sequences, this suggests exactly the values $18,\dots,21$.
A proof needs explicit error terms: an error of $1$ in $K_n$ would allow errors of up to $4$ in $a(n)$.

We prove the conjecture (Theorem~\ref{thm:main}) without using Hlawka's constant or any limit.
Two one-step inequalities (Lemmas~\ref{lem:up} and~\ref{lem:low}) show that $E(k)=S(k)-2\sqrt k$ is nonincreasing and $E(k)-2/\sqrt k$ is nondecreasing; on a window of three consecutive turns this gives $|a(n)-2\pi^2|<2.13$ once $n\ge71$ (Proposition~\ref{prop:core}).
The terms with $n\le70$ come from an interval computation evaluated by the Lean kernel.

\subsection*{Status and prior work}
The math0 literature referees found no prior proof and no counterexample.
A172164 still states the conjecture at its current revision (\#25, August~11, 2026), the comments in A137515 and A072895 are numerical, and W.~Lang's related entries (A296179, A296181, A295338, A295339) do not address the bound.
Searches of arXiv, MathOverflow, Mathematics Stack Exchange, the SeqFan list and the literature on the spiral \cite{Hl80,Da93,Br12} found no statement of the bound; Kociemba's web page on the spiral, linked from A072895, could only be seen through search snippets, which mention no proof.
The conjecture is not among the statements of the formal-conjectures repository \cite{FC} at the commit pinned by math0 (\texttt{c77fdd4}).
We therefore classify Theorem~\ref{thm:main} as a new result. The mathematics is modest, an effective version of Hlawka's formula (Remark~\ref{rem:hlawka}), and the bound may be known to experts without having been written down.

\section{Statement}\label{sec:statement}

Let $S$ be as above. For $n\ge1$ let
\[
  T(n)=\{\,i\ge1:\ 2\pi(n-1)\le S(i-1)\ \text{and}\ S(i)\le2\pi n\,\},\qquad c(n)=\#T(n),
\]
the set and number of triangles lying entirely in the $n$-th turn, and for $n\ge2$ let $a(n)=c(n)-c(n-1)$.

\begin{theorem}\label{thm:main}
For every integer $n\ge2$, $a(n)\in\{18,19,20,21\}$.
\end{theorem}

All four values occur: $a(2)=20$, $a(3)=19$, $a(7)=21$ and $a(8)=18$.

\subsection*{The OEIS conjecture}
The entry A172164 \cite{A172164} is named ``Differences between numbers of triangles entirely contained in two consecutive turns of Pythagoras's snail (Theodorus spiral).'' It has offset $2$, lists $69$ terms $a(2),\dots,a(70)$, and gives the example ``In the first turn, 16 triangles are complete. In the 2nd turn, there are 36 triangles completely included. The difference is 20.''
Its COMMENTS section contains, verbatim, the following conjecture by S\'ebastien Dumortier, dated January~27, 2010:
\begin{quote}\small
\texttt{Conjecture : The terms are only 18,19,20,21 (From the first thousand turns, there are 2,3\% of 18, 36,5\% of 19, 46,2\% of 20 and 15\% of 21). No period found. Probably due to Pi transcendence.}
\end{quote}

\begin{conjecture}[Dumortier]\label{conj}
Every term $a(n)$, $n\ge2$, of A172164 is one of $18,19,20,21$.
\end{conjecture}

The percentages and the last two sentences are observations, not part of the formalized claim.

\subsection*{Lean statement}
The statement is the definition \decl{Math0.Statements.A172164.statement} in the math0 statement file \path{lean/Math0/Statements/OEIS/A172164.lean}, which was reviewed for fidelity and frozen (SHA-256 \texttt{948ca46a\dots b30be5f4}) before any proof attempt.
The definitions and the statement read verbatim (docstrings, blank lines, the auxiliary lemmas and the data checks omitted):\par\nobreak\noindent
\begin{minipage}{\linewidth}
\begin{lstlisting}
namespace Math0.Statements.A172164
open Real Finset Filter Topology
noncomputable def S (k : ℕ) : ℝ := ∑ i ∈ range k, arctan (1 / √((i : ℝ) + 1))
def turnTriangles (n : ℕ) : Set ℕ :=
  {i | 1 ≤ i ∧ 2 * π * ((n : ℝ) - 1) ≤ S (i - 1) ∧ S i ≤ 2 * π * n}
noncomputable def c (n : ℕ) : ℕ := (turnTriangles n).ncard
noncomputable def a (n : ℕ) : ℤ := (c n : ℤ) - c (n - 1)
def statement : Prop := ∀ n : ℕ, 2 ≤ n → a n = 18 ∨ a n = 19 ∨ a n = 20 ∨ a n = 21
\end{lstlisting}
\end{minipage}

\emph{Encoding.} \texttt{S k} is $\sum_{i=0}^{k-1}\arctan(1/\sqrt{i+1})=S(k)$, and \texttt{turnTriangles n} is $T(n)$: \texttt{i - 1} is guarded by $1\le i$, and $n-1$ in the turn bound is computed in $\mathbb R$.
\texttt{Set.ncard} returns $0$ on infinite sets, but \decl{turnTriangles_finite} shows that $T(n)\subseteq[0,81n^2+2)$, using $S(k)\ge\frac56\sqrt k$ for $k\ge2$ (\decl{S_ge}), so \texttt{c n} is the true count.
\texttt{a n} is computed in $\mathbb Z$, and for $n\ge2$ the index $n-1\ge1$ is the previous turn.
The turn bounds are closed, the natural reading of ``entirely contained''; with open bounds triangle $1$, which starts at $S(0)=0$, would not count for turn $1$, contradicting the example $c(1)=16$. Otherwise the readings differ only if $S(k)\in2\pi\mathbb Z$ for some $k\ge1$. This does not happen for $k\le K_{70}$ (Proposition~\ref{prop:crossing}); we do not know whether it can happen at all, and Section~\ref{sec:proof} does not assume that it cannot.

\emph{Checks against the OEIS data.} The statement file proves, against the real-number definitions, that $c(1)=16$ and $c(2)=36$, as in the example (\decl{c_one_two}), and that $a(2),\dots,a(70)$ are the $69$ listed terms (\decl{a_data}).
Along the way it certifies the turn boundaries $K_1,\dots,K_{70}$, the first $42$ of which are the $42$ listed terms of A137515 (\decl{Kseq_A137515}).
These checks rest on the interval computation of Section~\ref{sec:small}, and the case $n\le70$ of Theorem~\ref{thm:main} uses them.

\section{Proof}\label{sec:proof}

Write $f(j)=\arctan(1/\sqrt j)$ and $E(k)=S(k)-2\sqrt k$.
Each lemma names its Lean declaration: unmarked names are private lemmas of the proof file, names marked (St) are in the statement file.
For $0\le x<1$ the terms of $\arctan x=\sum_{i\ge0}(-1)^ix^{2i+1}/(2i+1)$ decrease, so partial sums with an even (odd) number of terms are lower (upper) bounds (\decl{arctan_alt_bounds}, St); in particular $x-x^3/3\le\arctan x\le x$.

\subsection{Two one-step inequalities}

\begin{lemma}[\decl{step_upper}]\label{lem:up}
For every $k\ge0$, $f(k+1)\le2(\sqrt{k+1}-\sqrt k)$.
\end{lemma}

\begin{proof}
Let $p=\sqrt k$ and $q=\sqrt{k+1}\ge1$, so $q^2=p^2+1$.
If $q=1$ then $f(1)=\pi/4<1=1/q$; if $q>1$ then $\arctan(1/q)\le1/q$ by the alternating-series bound.
Finally $2q(q-p)-1=q^2+(q^2-1)-2pq=(q-p)^2\ge0$, so $1/q\le2(q-p)$.
\end{proof}

\begin{lemma}[\decl{step_lower}]\label{lem:low}
For every $k\ge1$, $f(k+1)\ge2(\sqrt{k+1}-\sqrt k)+2/\sqrt{k+1}-2/\sqrt k$.
\end{lemma}

\begin{proof}
Let $p=\sqrt k$, $q=\sqrt{k+1}$, so $0<p<q$ and $p^2=q^2-1$, and $x=1/q<1$.
By the alternating-series bound $f(k+1)\ge x-x^3/3$, and a direct computation gives
\[
  x-\frac{x^3}3-\Bigl(2(q-p)+\frac2q-\frac2p\Bigr)=\frac{N}{3pq^3},\qquad N=6q^3-3pq^2-p-6pq^3(q-p).
\]
Using $p^2=q^2-1$ one checks the identity $(q+p)N=3q^3(q-p)+q(2q-p)+1$, whose right-hand side is positive. Hence $N>0$.
\end{proof}

\begin{corollary}[\decl{E_anti}, \decl{E_mono}]\label{cor:E}
\textup{(a)} If $k\le m$ then $E(m)\le E(k)$; in particular $E(k)\le E(0)=0$ for all $k$.
\textup{(b)} If $1\le k\le m$ then $E(k)-2/\sqrt k\le E(m)-2/\sqrt m$.
\end{corollary}

\begin{proof}
Since $S(k+1)=S(k)+f(k+1)$ (\decl{S_succ}, St), Lemma~\ref{lem:up} says $E(k+1)\le E(k)$ and Lemma~\ref{lem:low} says $E(k)-2/\sqrt k\le E(k+1)-2/\sqrt{k+1}$ for $k\ge1$. Induct on $m$.
\end{proof}

\begin{remark}\label{rem:hlawka}
Corollary~\ref{cor:E} implies that $E(k)$ converges to a limit, Hlawka's constant $\kappa$ \cite{Hl80}, with $0\le E(k)-\kappa\le2/\sqrt k$ for $k\ge1$. This is not used below and not formalized.
\end{remark}

\subsection{Counting the triangles of a turn}
$S$ is strictly increasing with $S(0)=0$ (\decl{S_strictMono}, \decl{S_zero}, St), unbounded because $S(k)\ge\frac56\sqrt k$ (\decl{S_unbounded}), and $S(k+1)<S(k)+\pi/2$ because $\arctan<\pi/2$ (\decl{S_lt_succ_of_le}).

\begin{lemma}[\decl{exists_KL}]\label{lem:KL}
For every real $\theta\ge0$ there are $K,L\in\N$ such that $S(K)\le\theta<S(K+1)$; for all $j$, $\theta\le S(j)$ if and only if $L\le j$; and either $L=K$ and $S(K)=\theta$, or $L=K+1$.
\end{lemma}

\begin{proof}
Let $K+1$ be the least index with $\theta<S(K+1)$; it exists because $S$ is unbounded, and it is positive because $S(0)=0\le\theta$. Then $S(K)\le\theta$. Put $L=K$ if $S(K)=\theta$ and $L=K+1$ otherwise; the equivalence follows from strict monotonicity.
\end{proof}

\begin{lemma}[\decl{c_formula}]\label{lem:count}
If $S(K)\le2\pi m<S(K+1)$, $L\le K$, and $2\pi(m-1)\le S(j)\Leftrightarrow L\le j$ for all $j$, then $c(m)=K-L$.
\end{lemma}

\begin{proof}
For $i\ge1$: $2\pi(m-1)\le S(i-1)$ if and only if $i\ge L+1$, and $S(i)\le2\pi m$ if and only if $i\le K$. So $T(m)=\{L+1,\dots,K\}$, which has $K-L$ elements (\decl{ncard_Icc_nat}, St).
\end{proof}

\subsection{The squeeze on three consecutive turns}

\begin{lemma}[\decl{bound_X}]\label{lem:X}
Let $0<\pi'<3.1416$ and let $P,Q,V,T_0$ be reals with $200\le P\le Q$, $0\le V\le1/100$, $VP^2\le2(Q-P)$, $2P\le T_0<2P+1/P$ and $2Q\le T_0+4\pi'+V$. Then $X:=V(T_0+2\pi')+V^2/2\le0.13$.
\end{lemma}

\begin{proof}
From the hypotheses, $Q-P\le\frac12(T_0+4\pi'+V)-P<\frac1{2P}+2\pi'+\frac V2\le2\pi'+\frac3{400}$, and $T_1:=T_0+2\pi'<2P+2\pi'+\frac1{200}$. Hence
\[
  VT_1P^2\le2(Q-P)T_1\le2\Bigl(2\pi'+\tfrac3{400}\Bigr)\Bigl(2P+2\pi'+\tfrac1{200}\Bigr)\le2\cdot6.291\,(2P+6.291)\le0.129\,P^2,
\]
the last step because $0.129P^2-25.164P-79.16>0$ for $P\ge200$. So $VT_1\le0.129$, and $V^2/2\le0.001$.
\end{proof}

\begin{lemma}[\decl{final_bound}]\label{lem:final}
Let $3.14<\pi'<3.1416$, let $T_0,V,K_1,K_2,L_0,L_1$ be reals with $X=V(T_0+2\pi')+V^2/2\le0.13$, and put $T_1=T_0+2\pi'$, $T_2=T_0+4\pi'$. If
\begin{align*}
  &\text{(F1)}\ 4K_2\le(T_2+V)^2, &&\text{(F2)}\ T_2^2<4(K_2+1), &&\text{(F3)}\ 4L_0\le(T_0+V)^2+4,\\
  &\text{(F4)}\ T_0^2\le4L_0, &&\text{(F5)}\ T_1^2-2<2(K_1+L_1), &&\text{(F6)}\ 2(K_1+L_1)\le(T_1+V)^2+2,
\end{align*}
then $A=(K_2-L_1)-(K_1-L_0)$ satisfies $17<A<22$.
\end{lemma}

\begin{proof}
We have $4A=4K_2+4L_0-2\cdot2(K_1+L_1)$ and the identities
\[
  (T_2+V)^2+(T_0+V)^2-2T_1^2=8\pi'^2+4VT_1+2V^2,\qquad T_2^2+T_0^2-2(T_1+V)^2=8\pi'^2-4VT_1-2V^2 .
\]
By (F1), (F3), (F5), $4A<8\pi'^2+4X+8$; by (F2), (F4), (F6), $4A>8\pi'^2-4X-8$. So $|A-2\pi'^2|<2+X\le2.13$, and $9.85<\pi'^2<9.87$ gives $17<A<22$.
\end{proof}

\begin{proposition}[\decl{core}]\label{prop:core}
Let $r\ge69$ be real and $\theta_i=2\pi(r+i)$ for $i=0,1,2$. Let $K_0,K_1,K_2\in\N$ satisfy $S(K_i)\le\theta_i<S(K_i+1)$, and let $L_0,L_1\in\N$ satisfy, for $i=0,1$: either $L_i=K_i$ and $S(K_i)=\theta_i$, or $L_i=K_i+1$. Then $17<(K_2-L_1)-(K_1-L_0)<22$.
\end{proposition}

\begin{proof}
Since $S(K_0)\le\theta_0<\theta_1<S(K_1+1)$ and $S$ is increasing, $K_0\le K_1$; likewise $K_1\le K_2$.
Let $e=E(K_0)$, so $e\le0$ by Corollary~\ref{cor:E}(a), and put $T_i=\theta_i-e$; thus $T_1=T_0+2\pi$, $T_2=T_0+4\pi$.
By Corollary~\ref{cor:E}(a), $E(k)\le e$ for $k\ge K_0$, so:
(i) if $K\ge K_0$ and $\theta<S(K+1)$, then $\theta-e<2\sqrt{K+1}$; and (ii) if $K\ge K_0$ and $S(K)=\theta$, then $\theta-e\le2\sqrt K$.
By (i) with $K=K_0$ and $e\le0$, $2\sqrt{K_0+1}>2\pi r\ge2\cdot3.14\cdot69$, so $\sqrt{K_0+1}>216$ and $K_0\ge40000$.
Put $P=\sqrt{K_0}\ge200$, $Q=\sqrt{K_2}\ge P$ and $V=2/P-2/Q$. Then $V\ge0$, $V\le2/P\le1/100$, and $VPQ=2(Q-P)$ gives $VP^2\le2(Q-P)$.
For $K_0\le k\le K_2$, Corollary~\ref{cor:E}(b) (applicable as $K_0\ge1$) gives $E(k)\ge e-2/P+2/\sqrt k\ge e-V$, so:
(iii) if $K_0\le K\le K_2$ and $S(K)\le\theta$, then $2\sqrt K\le\theta-e+V$.
Next, $T_0=\theta_0-S(K_0)+2P\ge2P$; and since $(P+\frac1{2P})^2\ge K_0+1$, (i) gives $T_0<2\sqrt{K_0+1}\le2P+1/P$. By (iii) with $K=K_2$, $2Q\le T_2+V=T_0+4\pi+V$.
With $3.14<\pi<3.1416$ (Mathlib's \decl{Real.pi_gt_d2}, \decl{Real.pi_lt_d4}), Lemma~\ref{lem:X} gives $X\le0.13$.
It remains to verify (F1)--(F6) of Lemma~\ref{lem:final}; each follows by squaring inequalities between nonnegative numbers.
(F1) is (iii) at $K_2$, and (F2) is (i) at $K_2$.
For (F3), $L_0\le K_0+1$ and $4K_0=4P^2\le T_0^2\le(T_0+V)^2$.
For (F4): if $L_0=K_0$ and $S(K_0)=\theta_0$, then (ii) gives $T_0\le2\sqrt{K_0}$, so $T_0^2\le4L_0$; if $L_0=K_0+1$, (i) at $K_0$ gives $T_0^2<4L_0$.
For (F5): if $L_1=K_1$ and $S(K_1)=\theta_1$, (ii) gives $T_1^2\le4K_1=2(K_1+L_1)$; if $L_1=K_1+1$, (i) gives $T_1^2<4(K_1+1)=2(K_1+L_1)+2$.
For (F6), $L_1\le K_1+1$ and (iii) at $K_1$ gives $4K_1\le(T_1+V)^2$.
Lemma~\ref{lem:final} with $\pi'=\pi$ completes the proof.
\end{proof}

\subsection{The turns \texorpdfstring{$n\le70$}{n <= 70}}\label{sec:small}
All declarations named in this subsection are in the statement file.
Let $K_1,\dots,K_{70}$ be the numbers $16,53,109,\dots,47456,48835$ of the list \decl{Kseq}.

\begin{proposition}[\decl{crossing}, St]\label{prop:crossing}
For $1\le n\le70$, $S(K_n)<2\pi n<S(K_n+1)$.
\end{proposition}

\begin{proof}[The certificate]
Write $S(k+1)=\pi/4+R(k)$ with $R(k)=\sum_{j=2}^{k+1}f(j)$ (\decl{S_succ_eq_R}).
For $j\ge2$ and $x=1/\sqrt j<1$ we have $x^2=1/j$, and the alternating-series bounds with $2r$ and $2r+1$ terms read $x(A_r-B_r)\le f(j)\le x(A_{r+1}-B_r)$, where $A_r=\sum_{i<r}j^{-2i}/(4i+1)$ and $B_r=\sum_{i<r}j^{-2i-1}/(4i+3)$ (\decl{alt_sum_even}, \decl{alt_sum_odd}).
Let $D=2^{50}$. If $q_{\rm lo},q_{\rm hi}\in\N$ satisfy $q_{\rm lo}^2j\le2^{100}\le q_{\rm hi}^2j$, then $q_{\rm lo}\le Dx\le q_{\rm hi}$, and rounding every term down to an integer loses less than $1$ per term (\decl{term_lower}, \decl{term_upper}):
\begin{align*}
  \textstyle\sum_{i<r}\bigl\lfloor\frac{q_{\rm lo}}{(4i+1)j^{2i}}\bigr\rfloor-\sum_{i<r}\bigl\lfloor\frac{q_{\rm lo}}{(4i+3)j^{2i+1}}\bigr\rfloor-r
  &\le D\,f(j)\\
  &\le\textstyle\sum_{i\le r}\bigl\lfloor\frac{q_{\rm hi}}{(4i+1)j^{2i}}\bigr\rfloor+r+1-\sum_{i<r}\bigl\lfloor\frac{q_{\rm hi}}{(4i+3)j^{2i+1}}\bigr\rfloor .
\end{align*}
A program written with recursors and primitive operations on natural numbers (\decl{walkK}) adds these bounds for $j=2,3,\dots$ into running sums $a,b,c,d$ with the invariant $a-b\le D\,R(k)\le c-d$, proved for every input (\decl{walkK_spec}).
It finds $q_{\rm lo},q_{\rm hi}$ by a Newton iteration and checks $q_{\rm lo}^2j\le2^{100}\le q_{\rm hi}^2j$, falling back to the valid bounds $0$ and $D$ if a check fails; the number $r$ of terms affects only the precision.
At each target $(n,K_n)$ the program \decl{checkK} tests, in integer arithmetic, that $(c-d)/D<\pi_{\rm lo}(8n-1)/4$ for $k=K_n-1$ and $(a-b)/D>\pi_{\rm hi}(8n-1)/4$ for $k=K_n$, where $\pi_{\rm lo}=3.14159265358979323846<\pi<\pi_{\rm hi}=\pi_{\rm lo}+10^{-20}$ (Mathlib's \decl{Real.pi_gt_d20}, \decl{Real.pi_lt_d20}).
As $2\pi n-\pi/4=\pi(8n-1)/4$, the tests give $S(K_n)<2\pi n<S(K_n+1)$, and \decl{checkK_sound} states that a run returning \texttt{true} certifies all targets.
The theorem \decl{check_targets}, that the run on the $70$ targets returns \texttt{true}, is proved by \texttt{decide +kernel}: the Lean kernel evaluates the program over $j=2,\dots,48836$ with its built-in arbitrary-precision arithmetic on natural number literals.
The accumulated enclosure width is about $1.4\cdot10^{-7}$, while the smallest of the gaps $2\pi n-S(K_n)$ and $S(K_n+1)-2\pi n$ is about $2.1\cdot10^{-5}$ (at $n=53$).
\end{proof}

By Proposition~\ref{prop:crossing} and strict monotonicity, $T(1)=\{1,\dots,K_1\}$ and $T(n)=\{K_{n-1}+2,\dots,K_n\}$ for $2\le n\le70$ (\decl{turnTriangles_one}, \decl{turnTriangles_eq}), so $c(1)=K_1$ and $c(n)=K_n-K_{n-1}-1$ (\decl{c_eq}), in accordance with the entry's formula ``second forward difference of A072895''. Hence \decl{a_data}: $a(2),\dots,a(70)$ are the $69$ listed terms, all in $\{18,19,20,21\}$.

\subsection{Proof of Theorem~\ref{thm:main} (\texorpdfstring{\texttt{result}}{result})}
Let $n\ge2$. If $n\le70$, $a(n)$ occurs in the list of \decl{a_data}, and every entry of that list is $18$, $19$, $20$ or $21$ (checked by \texttt{decide}).
Let $n\ge71$ and $r=n-2\ge69$. Apply Lemma~\ref{lem:KL} to $\theta_i=2\pi(r+i)$, $i=0,1,2$, to get $K_i,L_i$.
Since $S(K_0+1)<\theta_0+\pi/2<\theta_1<S(K_1+1)$, we have $L_0\le K_0+1\le K_1$, and likewise $L_1\le K_2$.
Turn $n$ is $[\theta_1,\theta_2]$ and turn $n-1$ is $[\theta_0,\theta_1]$, so Lemma~\ref{lem:count} gives $c(n)=K_2-L_1$ and $c(n-1)=K_1-L_0$.
By Proposition~\ref{prop:core}, the integer $a(n)=(K_2-L_1)-(K_1-L_0)$ satisfies $17<a(n)<22$. \qed

\section{Machine verification}\label{sec:lean}

\subsection*{Files}
The proof is in \path{lean/Math0/Frontier/OEIS/A172164.lean} of the math0 repository, \url{https://github.com/Foundation-0/math0}; it imports the frozen statement file and the audit commands of \path{lean/Math0/Audit.lean}.
The main theorem is \decl{Math0.Frontier.A172164.result}, of type \decl{statement}, preceded by the private lemmas named in Section~\ref{sec:proof}; the case $n\le70$ uses \decl{a_data} of the statement file.
Neither file uses \texttt{sorry}, \texttt{native\_decide}, \texttt{set\_option}, custom axioms or custom syntax.

\subsection*{Audit}
The file ends with the command \texttt{\#math0\_settles} applied to \decl{result} and \decl{statement}.
This command (\path{lean/Math0/Audit.lean}) fails unless \decl{statement} is a \texttt{Prop}-valued definition, \decl{result} has syntactically the type \decl{statement} (or its negation), and both depend on no axioms besides \decl{propext}, \decl{Classical.choice} and \decl{Quot.sound}; this rejects \decl{sorryAx}, the axiom \decl{Lean.ofReduceBool} behind \texttt{native\_decide}, and custom axioms.
Since \decl{result} depends on \decl{a_data} and hence on \decl{check_targets}, the axiom check also covers the interval computation.
Running \texttt{lake env lean Math0/Frontier/OEIS/A172164.lean} in the directory \path{lean} prints\par\nobreak\noindent
\begin{minipage}{\linewidth}
\begin{lstlisting}[basicstyle=\ttfamily\scriptsize]
math0 audit: 'Math0.Statements.A172164.statement' depends on axioms [Classical.choice, Quot.sound, propext]
math0 audit: 'Math0.Frontier.A172164.result' depends on axioms [Classical.choice, Quot.sound, propext]
math0 statement check: 'Math0.Frontier.A172164.result' proves 'Math0.Statements.A172164.statement'
\end{lstlisting}
\end{minipage}

\subsection*{Toolchain and reproduction}
We use Lean~4 \cite{Lean4} \texttt{v4.33.1} and Mathlib \cite{Mathlib} \texttt{v4.33.1}, pinned in \path{lean/}; the repository also pins formal-conjectures \cite{FC} at commit \texttt{c77fdd4}, which this result does not use.
To reproduce, run \path{scripts/verify.sh} from the repository root: it fetches the Mathlib build cache, builds all math0 results (re-running every audit command) and checks that the pinned dependencies are unmodified.
To check this result alone, run \texttt{lake build Math0.Statements.OEIS.A172164} (about a minute, half of it the kernel evaluation of \decl{check_targets}) and then \texttt{lake env lean Math0/Frontier/OEIS/A172164.lean} (about half a minute) in the directory \path{lean}.

{\footnotesize
\begin{thebibliography}{99}

\bibitem{Br12}
D.~Brink, \emph{The spiral of Theodorus and sums of zeta-values at the half-integers}, Amer. Math. Monthly \textbf{119} (2012), no.~9, 779--786.

\bibitem{Da93}
P.~J.~Davis, \emph{Spirals: From Theodorus to Chaos}, A K Peters, Wellesley, MA, 1993.

\bibitem{FC}
M.~Firsching, P.~Lezeau, S.~Mercuri, M.~Z. Horv\'ath, Y.~Dillies, C.~S\"onne, E.~Wieser, F.~Zhang, T.~Hubert, B.~Ag\"uera y Arcas, and P.~Kohli, \emph{Formal Conjectures: An Open and Evolving Benchmark for Verified Discovery in Mathematics}, arXiv:2605.13171, 2026. Repository: \url{https://github.com/google-deepmind/formal-conjectures}.

\bibitem{Hl80}
E.~Hlawka, \emph{Gleichverteilung und Quadratwurzelschnecke}, Monatsh. Math. \textbf{89} (1980), 19--44.

\bibitem{Lean4}
L.~de Moura and S.~Ullrich, \emph{The Lean 4 theorem prover and programming language}, in: Automated Deduction -- CADE 28, Lecture Notes in Comput. Sci. \textbf{12699}, Springer, 2021, 625--635.

\bibitem{Mathlib}
The mathlib Community, \emph{The Lean mathematical library}, in: Proceedings of the 9th ACM SIGPLAN International Conference on Certified Programs and Proofs (CPP 2020), ACM, 2020, 367--381.

\bibitem{A072895}
The OEIS Foundation Inc., \url{https://oeis.org/A072895} (least $k$ for the Theodorus spiral to complete $n$ revolutions), entry by R.~G.~Wilson~v, Jul.~29, 2002.

\bibitem{A105459}
The OEIS Foundation Inc., \url{https://oeis.org/A105459} (decimal expansion of Hlawka's Schneckenkonstante), entry by D.~Brink, Jun.~13, 2011.

\bibitem{A137515}
The OEIS Foundation Inc., \url{https://oeis.org/A137515} (maximal number of right triangles in $n$ turns of Pythagoras's snail), entry by S.~Dumortier, Apr.~23, 2008.

\bibitem{A172164}
The OEIS Foundation Inc., \url{https://oeis.org/A172164} (differences between numbers of triangles entirely contained in two consecutive turns of Pythagoras's snail), entry by S.~Dumortier, Jan.~27, 2010; revision \#25, Aug.~11, 2026.

\end{thebibliography}
}

\end{document}
